\section{Occupation of retained prefixes}
\label{sec:poisson_word_occupation}

Let $I$ be a finite alphabet, let $S=\{i\in I:p(i)\}$, and let
$R_p(w)$ denote the retained coordinate of $P_p(w)$ from
Definition~\ref{def:poisson_word_partition}.  Write $\mathcal W_S$
for the words over $S$ and $\mu_{S,s}$ for their Poissonized law.
At a marked position $j$ of a full word $w$, the observable below uses
$R_p(w_{<j})$: the retained letters strictly before that position.
Thus each contribution uses its own retained prefix.

\begin{theorem}[Retained-prefix occupation identity]
\label{thm:poisson_word_partition_prefix}
\lean{PoissonWord.lintegral_sum_partition_prefix}
\leanok
\uses{def:poisson_word_partition, def:poisson_word_split}
For $T\geq0$ and any $F:\mathcal W_S\times I\to[0,\infty]$,
\begin{equation}
\label{eq:poisson_word_partition_prefix}
\begin{aligned}
 &\int_{\mathcal W}\sum_{j=0}^{|w|-1}
 F(R_p(w_{<j}),w_j)\,d\mu_{I,T}(w)\\
 &\qquad=\int_0^T\sum_{i\in I}\int_{\mathcal W_S}
 F(u,i)\,d\mu_{S,s}(u)\,ds.
\end{aligned}
\end{equation}
These are nonnegative integrals; both sides may be infinite.
\end{theorem}
\begin{proof}
\leanok
\uses{cor:poisson_word_insertion_prefix, thm:poisson_word_partition}
Apply the prefix insertion identity to $G(u,i)=F(R_p(u),i)$.
For each insertion time $s$, the retained marginal of a word with law
$\mu_{I,s}$ is $\mu_{S,s}$ by subset thinning.  Transporting the inner
integral along $R_p$ gives the asserted formula.  No independence
assumption about the different random prefixes is needed.
\end{proof}

\begin{corollary}[Occupation at omitted letters]
\label{cor:poisson_word_omitted_prefix}
\lean{PoissonWord.lintegral_sum_omitted_partition_prefix}
\leanok
\uses{thm:poisson_word_partition_prefix}
With the same assumptions,
\begin{equation}
\label{eq:poisson_word_omitted_prefix}
\begin{aligned}
 &\int_{\mathcal W}
 \sum_{\substack{0\leq j<|w|\\w_j\notin S}}
 F(R_p(w_{<j}),w_j)\,d\mu_{I,T}(w)\\
 &\qquad=\int_0^T\sum_{i\in I\setminus S}
 \int_{\mathcal W_S}F(u,i)\,d\mu_{S,s}(u)\,ds.
\end{aligned}
\end{equation}
\end{corollary}
\begin{proof}
\leanok
Replace $F(u,i)$ by zero for $i\in S$ and leave it unchanged otherwise
in~(\ref{eq:poisson_word_partition_prefix}).
\end{proof}

\begin{corollary}[Expected omitted count]
\label{cor:poisson_word_omitted_count}
\lean{PoissonWord.lintegral_omitted_count}
\leanok
\uses{cor:poisson_word_omitted_prefix}
The expected number of omitted letters is
\begin{align}
 \int_{\mathcal W}\#\{j<|w|:w_j\notin S\}\,d\mu_{I,T}(w)
 &=|I\setminus S|\,T.
 \label{eq:poisson_word_omitted_count}
\end{align}
\end{corollary}
\begin{proof}
\leanok
\uses{lem:poisson_word_normalization}
Set $F=1$ in~(\ref{eq:poisson_word_omitted_prefix}).
Each retained-word measure has mass one, leaving $|I\setminus S|$
integrated over the time interval.
\end{proof}

All three identities include $T=0$, $S=\varnothing$, and $S=I$.
They supply the probability integration for the earlier retained factors
in \cite{OpenAI2026AreaLaw}, \texttt{09-amplification.tex}, lines 203--209.
The operator telescope, clock paths, propagation estimates, and final
spatial-omission estimate remain separate inputs.
% Source revision: adc7f1241b42e322a6451854ab7e4b4c146bf78a.
